% !TEX root = ../main.tex
% Figure: PWEM plane wave expansion method workflow
% Chapter: ch12 - Plane Wave Expansion Method and Eigenvalue Solvers
% Description: Flowchart showing the computational steps in PWEM calculation

\begin{figure}[htbp]
  \centering
  \resizebox{\textwidth}{!}{%
  \begin{tikzpicture}[scale=0.92, transform shape,
    block/.style={rectangle, draw, rounded corners, fill=#1!20, minimum width=3cm, minimum height=0.8cm, align=center},
    decision/.style={diamond, draw, aspect=2, fill=#1!20, minimum width=2.5cm, align=center},
    arrow/.style={->, >=stealth, thick}
  ]
    % Input stage
    \node[block=blue, minimum height=1.2cm] (input) at (0,8) {输入结构参数\\晶格类型$a$,填充因子$f$,折射率$n$};
    
    % Geometry setup
    \node[block=green] (geometry) at (0,6.5) {构建介电函数$\varepsilon(\mathbf{G})$\\傅里叶分量计算};
    \draw[arrow] (input) -- (geometry);
    
    % Reciprocal lattice
    \node[block=yellow, right=2cm of geometry] (recip) {生成倒格矢$\{\mathbf{G}_n\}$\\截断至$N_{\mathrm{pw}}$个平面波};
    \draw[arrow] (geometry) -- (recip);
    
    % Master equation assembly
    \node[block=orange] (matrix) at (0,5) {组装本征方程\\$\sum_{\mathbf{G}'}\Theta_{\mathbf{G},\mathbf{G}'}H_{\mathbf{G}'}=\left(\frac{\omega}{c}\right)^2 H_{\mathbf{G}}$};
    \draw[arrow] (recip) |- (matrix);
    \draw[arrow] (geometry) -- (matrix);
    
    % K-point sampling
    \node[block=cyan, left=2cm of matrix] (kspace) {$\mathbf{k}$空间采样\\沿不可约路径取点};
    \draw[arrow] (kspace) -- (matrix);
    
    % Eigenvalue solver
    \node[decision=red] (solver) at (0,3.5) {求解大规模稀疏\\矩阵本征值？};
    \draw[arrow] (matrix) -- (solver);
    
    % Iterative refinement loop
    \node[block=purple] (iterate) at (-3,2) {增加$N_{\mathrm{pw}}$\\提高精度};
    \draw[arrow] (solver) -| node[pos=0.3,above] {否/收敛慢} (iterate);
    \draw[arrow] (iterate) |- (matrix);
    
    % Post-processing
    \node[block=pink] (postproc) at (3,2) {后处理\\群速度$v_g=d\omega/dk$\\态密度 DOS};
    \draw[arrow] (solver) -| node[pos=0.3,above] {是} (postproc);
    
    % Output
    \node[block=blue, minimum height=1.2cm] (output) at (0,0.5) {输出结果\\能带图$\omega_n(\mathbf{k})$, 场分布，偏振信息};
    \draw[arrow] (postproc) |- (output);
    \draw[arrow] (iterate) |- (output);
    
    % Computational cost annotation
    \node[draw, dashed, align=center, font=\small] at (-5,4) {计算复杂度:\\$O(N_{\mathrm{pw}}^3)$\\内存:$O(N_{\mathrm{pw}}^2)$};
    
    % Tips box
    \node[align=left, font=\footnotesize, anchor=south west] at (4.5,0.5) {
      \textbf{实用建议}:\\
      $\bullet$ $N_{\mathrm{pw}} \geq 441$ ($21\times21$) 保证收敛\\
      $\bullet$ 使用 ARPACK 迭代求解器加速\\
      $\bullet$ 对称性简化可减少 8 倍计算量
    };
  \end{tikzpicture}%
  }
  \caption{PWEM 的标准计算流程：由几何与材料参数构造倒格矢和介电函数傅里叶分量，离散 Maxwell 本征方程并沿高对称路径求解本征频率与本征矢量；若结果未收敛，则增大 $N_{\mathrm{pw}}$ 重新计算。输出的能带、模场和群速度等量可用于 PCSEL 的初始选模与参数扫描。}
  \label{fig:ch12_pwem_workflow}
\end{figure}
